\notechapter{N-20221210}{Binomial Theorem, Absolute-Value Estimates, and Set Exercises}

\section{Binomial theorem}

The first page begins with the binomial theorem:
\[
  (a+b)^n=\sum_{i=0}^n \binom ni a^i b^{n-i}.
\]
The handwritten derivation emphasizes selection.  When expanding
\[
  (a+b)(a+b)\cdots(a+b),
\]
one chooses either $a$ or $b$ from each factor.  A term $a^ib^{n-i}$ occurs when $a$ is chosen from exactly $i$ of the $n$ factors, and this can be done in $\binom ni$ ways.

The same idea proves the recurrence
\[
  \binom{n+1}{i}=\binom ni+\binom n{i-1},
\]
because in choosing $i$ objects from $n+1$, either the new object is not chosen or it is chosen.

\section{Absolute value of binomial sums}

The note then estimates expressions such as
\[
  |a+b|^n.
\]
Using the triangle inequality and the binomial theorem,
\[
  |a+b|^n
  \le (|a|+|b|)^n
  =\sum_{i=0}^n \binom ni |a|^i|b|^{n-i}.
\]
The important habit is to move absolute values inside after expanding only when one has a legitimate inequality.  The handwritten page marks the triangle inequality as the justification.

\section{A summation exercise}

One exercise asks for a sum of the form
\[
  \sum_{i=0}^n \binom ni.
\]
Putting $a=b=1$ in the binomial theorem gives
\[
  2^n=\sum_{i=0}^n \binom ni.
\]
Similarly, putting $a=1,b=-1$ gives
\[
  0=(1-1)^n=\sum_{i=0}^n (-1)^{n-i}\binom ni,
\]
which shows that the sum of even-indexed binomial coefficients equals the sum of odd-indexed ones when $n>0$.

\section{Set exercises}

The printed pages list elementary set exercises.  The notes use standard notation:
\[
  A\cup B,\qquad A\cap B,\qquad A\setminus B,\qquad A^c.
\]
One recurring task is to translate a verbal condition into set operations.  For example,
\[
  A\setminus(B\cup C)=(A\setminus B)\cap(A\setminus C).
\]
This is proved elementwise:
\[
  x\in A\setminus(B\cup C)
  \Longleftrightarrow
  x\in A,\ x\notin B,\ x\notin C.
\]

The exercises also ask for necessary and sufficient conditions for equalities such as
\[
  A\cup B=A,\qquad A\cap B=B.
\]
Both are equivalent to
\[
  B\subseteq A.
\]

\section{Counting subsets with properties}

For a finite set $X$ with $|X|=n$, the number of subsets is
\[
  2^n.
\]
If one imposes parity conditions, the count is often half of the total.  For example, the number of subsets of an $n$-element set with even cardinality is
\[
  2^{n-1},
\]
when $n\ge1$.  This follows from
\[
  \sum_{k\text{ even}}\binom nk
  =
  \sum_{k\text{ odd}}\binom nk
  =2^{n-1}.
\]

The printed exercises include conditions involving positive and negative numbers, odd and even elements, and sums of elements.  The general method is to split the set into independent categories and multiply the numbers of choices.

\section{Set-builder notation}

The note also works with sets such as
\[
  U=\{x:x=2n,\ n\in\mathbb Z\},\qquad
  V=\{x:x=3n,\ n\in\mathbb Z\}.
\]
Here $U$ is the set of even integers and $V$ is the set of multiples of $3$.  Then
\[
  U\cap V=\{x:x=6n,\ n\in\mathbb Z\}
\]
and
\[
  U\cup V
\]
is the set of integers divisible by $2$ or by $3$.

The point is to read set-builder notation as a logical predicate.

\section{Expanded synthesis and advanced context}

The note links binomial coefficients with set counting.  The binomial theorem counts choices from factors; subset counting counts choices from elements.  The same coefficients $\binom nk$ appear in both languages.  Set identities should be proved elementwise, while counting identities should be proved by telling a story about what is being chosen.



\paragraph{Expanded synthesis.}
The elementary geometric content of this chapter should be read through transformations and invariants. The earlier sections (`Binomial theorem`, `Absolute value of binomial sums`, `A summation exercise`, `Set exercises`, `Counting subsets with properties`) may use coordinates, angles, determinants, parametrizations, or integral formulas, but the organizing question is: which quantities survive the transformations naturally associated with the problem? Euclidean geometry preserves distances and angles; affine geometry preserves lines and ratios on a line; projective geometry preserves incidence and cross ratio; differential geometry preserves coordinate-independent objects such as forms, metrics, and orientations.

This perspective separates different layers of a problem. A dot product belongs to metric geometry. A determinant belongs to oriented volume. A cross ratio belongs to projective geometry. A line or surface integral of a differential form belongs to oriented topology. A scalar area integral belongs to the metric-induced measure. Confusing these layers is the source of many errors, especially in vector calculus where $dS$, $F\cdot n\,dS$, $dx\wedge dy$, and boundary orientation look similar but encode different structures.

When returning to the original examples, identify the natural object before computing. If the problem is about concurrence or collinearity, projective or affine tools may be natural. If it is about distance, angle, or orthogonality, metric tools are essential. If it is about flux or circulation, differential forms and orientation explain the signs. The advanced viewpoint makes the elementary solution more disciplined rather than more abstract for its own sake.


\section{Elementary-to-advanced bridge: binomial transform}

The binomial theorem gives
\[
  (1+x)^n=\sum_{k=0}^n \binom nk x^k.
\]
The binomial transform of a sequence $(a_k)$ is
\[
  b_n=\sum_{k=0}^n \binom nk a_k.
\]
It has an inverse:
\[
  a_n=\sum_{k=0}^n (-1)^{n-k}\binom nk b_k.
\]
This is Möbius inversion on the subset lattice in sequence form.

Thus binomial identities in the note can often be read as transformations between two coordinate systems for sequences.

\section{Elementary-to-advanced bridge: exponential generating functions and labelled structures}

For labelled combinatorial structures, the exponential generating function is
\[
  A(x)=\sum_{n\ge0}a_n\frac{x^n}{n!}.
\]
The factor $n!$ compensates for labellings.  Products of exponential generating functions correspond to splitting a labelled set into labelled parts.

For example, the exponential formula says that if a structure is a set of connected components, then its exponential generating function is the exponential of the generating function for connected components:
\[
  A(x)=\exp(C(x)).
\]
This is the advanced version of the elementary idea that a set can be decomposed into independent pieces.

\section{Elementary-to-advanced bridge: combinatorial species}

A combinatorial species is a functor from finite sets with bijections to finite sets of structures.  For each finite label set $U$, a species $F$ assigns a set $F[U]$ of structures on $U$; every bijection $U\to V$ transports structures from $F[U]$ to $F[V]$.

This functorial language records not only how many objects exist, but how they behave under relabelling.  The binomial theorem corresponds to the species identity for subsets: choosing a subset of an $n$-set means splitting labels into chosen and unchosen parts.

\section[Probability view of binomial coefficients]{Elementary-to-advanced bridge: probability distributions behind binomial coefficients}

The binomial coefficients also appear in probability.  If $X\sim \operatorname{Bin}(n,p)$, then
\[
  \mathbb P(X=k)=\binom nk p^k(1-p)^{n-k}.
\]
The binomial theorem ensures total probability equals $1$:
\[
  \sum_{k=0}^n\binom nk p^k(1-p)^{n-k}=1.
\]
The generating function is
\[
  \mathbb E[t^X]=(1-p+pt)^n.
\]
Thus binomial algebra, subset counting, and sums of independent Bernoulli random variables are the same structure.

% Second-pass deepening for waves 2-4: wave 3 A-C part 1
\section{Second-pass deepening: binomial identities as change of basis}

The polynomials
\[
  \binom x0,\ \binom x1,\ \binom x2,\ldots
\]
form a natural basis for polynomial functions on integers.  The finite difference operator acts simply:
\[
  \Delta \binom xk=\binom x{k-1}.
\]
Therefore many binomial identities are not coincidences; they are coordinate transformations in this basis.

For example, Vandermonde's identity
\[
  \sum_k \binom rk\binom s{n-k}=\binom{r+s}n
\]
is coefficient extraction from
\[
  (1+x)^r(1+x)^s=(1+x)^{r+s}.
\]
It also counts choosing $n$ elements from a disjoint union of sets of sizes $r$ and $s$.

\section{Second-pass deepening: species and the meaning of labels}

In labelled counting, the labels matter only up to bijection.  A species formalizes this by requiring every bijection of label sets to transport structures.  This turns counting into functorial combinatorics.

For instance, the species of subsets assigns to a finite set $U$ the set $\mathcal P(U)$.  A bijection $U\to V$ sends a subset of $U$ to the corresponding subset of $V$.  The binomial theorem is the decategorified shadow of this functorial behavior.
